% ============================================================
% main.tex — the article itself
%
% This file is structure and content. Every typographic
% decision lives in ffvstyle.sty, which also loads babel.
%
% Compile with LuaLaTeX, not pdflatex:
%   latexmk -lualatex main
% ============================================================
\documentclass[a4paper,10pt,twocolumn]{article}

% ============================================================
% ARTICLE STYLE
% ============================================================
\usepackage{ffvstyle}
% Uncomment for links with no colour at all:
% \hypersetup{hidelinks}

% ============================================================
% MATHEMATICS
%
% No amssymb: unicode-math supplies the symbols, and
% loading both breaks the build (\eth already defined).
% No bm either — bold maths is \symbf / \symbfit /
% \symbfup. mathtools is not loaded; the README lists the
% unicode-math-native replacements.
% ============================================================
\usepackage{amsmath}

% ============================================================
% CUSTOM PACKAGES AND COMMANDS
%
% Add additional packages and custom commands here.
% Examples that are often useful are commented out
% below — remove the % to activate them.
% ============================================================

% --- Additional mathematics packages ---
% \usepackage{slashed}        % Feynman slash for Dirac matrices
% \usepackage{tensor}         % Nice tensor indices
% \usepackage{physics2}       % then \usephysicsmodule{braket,ab.legacy}
%                             % for \ket, \bra, \abs, \norm
% \usepackage{cancel}         % Cancelled mathematics

% --- Additional text and layout packages ---
% \usepackage{enumitem}       % Fine control over lists
% \usepackage{siunitx}        % SI units: \qty{9.81}{\metre\per\second\squared}
% \usepackage{epigraph}       % Epigraphs at section starts

% --- Custom commands ---
% \newcommand{\R}{\mathbb{R}}
% \newcommand{\C}{\mathbb{C}}
% \newcommand{\mean}[1]{\langle #1 \rangle}

% --- Custom operators ---
% \DeclareMathOperator{\Tr}{Tr}
% \DeclareMathOperator{\sgn}{sgn}

% ============================================================
% METADATA — fill in here
% ============================================================
\title{ARTICLE TITLE\\[0.4em]
{\large Optional subtitle}}
\author{First Last}
\date{\today}

% ============================================================
% DOCUMENT
% ============================================================
\begin{document}

\twocolumn[
  \maketitle
  \begin{lead}
    Write the lead paragraph here. The lead is a short, self-contained
    paragraph in bold that gives the reader a quick overview of what
    the article is about, what the central argument is, and why it
    is interesting. Two to four sentences is usually about right.
    Citations are allowed here too~\cite{example2024}.
  \end{lead}
]

% ============================================================
\section{Introduction}

Write the introduction here. State the problem, give the necessary
background, and end with a short overview of the article's structure
if useful. Paragraphs are marked with indentation --- not with extra
line breaks.

% ============================================================
\section{Section Two}

Mathematics is set inline as $E = mc^2$ or as numbered equations:
\begin{equation}
  \hat{H}\psi = E\psi.
  \label{eq:example}
\end{equation}
Refer to equations with \eqref{eq:example} and to references with
\cite{example2024}.

% -------------------------------------------------
% FIGURE — single column (floats within one column)
% -------------------------------------------------
\begin{figure}[t]
  \centering
  % Replace with your own figure:
  % \includegraphics[width=\linewidth]{filename}
  \rule{0.9\linewidth}{3cm}     % placeholder
  \caption{This is a caption for a single-column figure.
           The text wraps automatically and is indented neatly
           below the label.}
  \label{fig:simple}
\end{figure}

Refer to figures in the text: figure~\ref{fig:simple} shows an
example. The placement argument \texttt{[t]} asks LaTeX to place
the figure at the top of a column --- the most common options are
\texttt{t} (top), \texttt{b} (bottom), \texttt{h} (here if possible),
and \texttt{p} (own page).

\subsection{Subsection}

Use subsections to structure longer sections.

% -------------------------------------------------
% FIGURE — two-column (spans both columns)
% -------------------------------------------------
\begin{figure*}[t]
  \centering
  % Replace with your own figure:
  % \includegraphics[width=0.9\linewidth]{filename}
  \rule{0.9\linewidth}{3.5cm}   % placeholder
  \caption{This is a two-column figure that spans the full page
           width. Use it for wide illustrations, multi-panel plots,
           or figures that need a lot of space to be readable.
           The \texttt{figure*} environment floats only to the top
           or bottom of pages.}
  \label{fig:wide}
\end{figure*}

% -------------------------------------------------
% TABLE
% -------------------------------------------------
\begin{table}[t]
  \centering
  \caption{An example of a tasteful table with \texttt{booktabs}.
           Vertical lines are avoided; horizontal lines are used
           sparingly (top, between header and data, bottom).}
  \label{tab:example}
  \begin{tabular}{lcr}
    \toprule
    Quantity & Symbol & Value \\
    \midrule
    Speed of light       & $c$      & $2.998 \times 10^8$ m/s \\
    Planck's constant    & $\hbar$  & $1.055 \times 10^{-34}$ Js \\
    Elementary charge    & $e$      & $1.602 \times 10^{-19}$ C \\
    \bottomrule
  \end{tabular}
\end{table}

% ============================================================
\section{Conclusion}

Summarize the main results and feel free to point toward open
questions or further work.

% ============================================================
% REFERENCES
% Compile order:
%   lualatex main
%   bibtex   main
%   lualatex main
%   lualatex main
% ============================================================
\bibliographystyle{unsrtnat}
\bibliography{references}

\end{document}
